\chapter*{插图索引}
\addcontentsline{toc}{chapter}{插图索引}
\begingroup
\small
\setlength{\tabcolsep}{4pt}
\renewcommand{\arraystretch}{1.15}
\newcommand{\FigureIndexLink}[2]{\hyperlink{page.#2}{#1}}
\newcommand{\FigureIndexRow}[3]{%
  \FigureIndexLink{#1}{#3} & \FigureIndexLink{#2}{#3} & \FigureIndexLink{#3}{#3} \\
}
\begin{longtable}{p{1.9cm}p{10.0cm}r}
\toprule
图号 & 简短标题 & 页码 \\
\midrule
\endfirsthead
\toprule
图号 & 简短标题 & 页码 \\
\midrule
\endhead
\FigureIndexRow{图1.1}{四种主要半导体激光器的结构对比}{6}
\FigureIndexRow{图1.2}{PCSEL 的纵向层状结构与面内反馈、法向辐射概念}{7}
\FigureIndexRow{图1.3}{光子晶体面发射激光器建模的三重视角与数据传递链}{11}
\FigureIndexRow{图2.1}{本书统一教学器件示意}{15}
\FigureIndexRow{图2.2}{本书采用的四层建模结构}{19}
\FigureIndexRow{图2.3}{全书逻辑依赖图}{23}
\FigureIndexRow{图2.4}{“六条红线”可视化速记图}{23}
\FigureIndexRow{图3.1}{复频率平面示意}{30}
\FigureIndexRow{图4.1}{正方晶格实空间原胞与第一布里渊区}{42}
\FigureIndexRow{图4.2}{二维实空间晶格与倒空间晶格示意}{42}
\FigureIndexRow{图4.3}{方格晶格第一布里渊区构造示意}{43}
\FigureIndexRow{图4.4}{方格晶格第一布里渊区高对称点路径}{43}
\FigureIndexRow{图4.5}{三角晶格第一布里渊区示意}{44}
\FigureIndexRow{图5.1}{Γ 点与非 Γ 点模式近场远场对比}{54}
\FigureIndexRow{图5.2}{正方晶格圆孔教学几何与路径示意}{58}
\FigureIndexRow{图5.3}{能带图读图示意}{59}
\FigureIndexRow{图6.1}{动量折叠与光锥耦合示意}{71}
\FigureIndexRow{图6.2}{平板导模色散与光锥边界示意}{71}
\FigureIndexRow{图6.3}{三类主要面发射机制对比}{72}
\FigureIndexRow{图7.1}{损耗机制分解示意图}{80}
\FigureIndexRow{图7.2}{阈值增益分解示意图}{84}
\FigureIndexRow{图8.1}{近场对称性到远场角谱映射示意}{96}
\FigureIndexRow{图8.2}{远场辐射方向图示例}{96}
\FigureIndexRow{图8.3}{偏振调控方法对比}{99}
\FigureIndexRow{图9.1}{远场偏振涡旋最小示意}{109}
\FigureIndexRow{图10.1}{最小四波 CWT 图像}{114}
\FigureIndexRow{图11.1}{耦合波升级路径示意}{141}
\FigureIndexRow{图12.1}{端口-腔模-端口最小模型}{150}
\FigureIndexRow{图12.2}{微扰诱导模式分裂示意}{156}
\FigureIndexRow{图13.1}{PWEM 标准计算流程图}{164}
\FigureIndexRow{图13.2}{平面波矩阵结构与截断示意}{165}
\FigureIndexRow{图13.3}{二维标量能带实算结果}{169}
\FigureIndexRow{图14.1}{RCWA 阶梯离散策略}{179}
\FigureIndexRow{图15.1}{三维 FDTD 仿真配置示意}{189}
\FigureIndexRow{图16.1}{数值方法性能雷达对比}{199}
\FigureIndexRow{图17.1}{典型外延堆栈截面示意}{212}
\FigureIndexRow{图17.2}{导带/价带台阶示意}{213}
\FigureIndexRow{图17.3}{外延层纵向结构示意}{214}
\FigureIndexRow{图17.4}{量子阱重叠与损耗重叠示意}{217}
\FigureIndexRow{图18.1}{压缩应变量子阱能带结构示意}{222}
\FigureIndexRow{图18.2}{多量子阱与阻挡层能带示意}{224}
\FigureIndexRow{图19.1}{横向扩展层与注入均匀性关系示意}{233}
\FigureIndexRow{图19.2}{电流拥挤效应示意}{237}
\FigureIndexRow{图20.1}{热效应多物理场耦合路径图}{244}
\FigureIndexRow{图20.2}{热效应多物理场耦合路径图（教学版）}{244}
\FigureIndexRow{图20.3}{热路等效电路类比图}{245}
\FigureIndexRow{图20.4}{热-电-光耦合闭环量传递示意}{247}
\FigureIndexRow{图21.1}{大面积阵列相干锁定机制对比}{252}
\FigureIndexRow{图22.1}{工艺误差传播图}{259}
\FigureIndexRow{图23.1}{动态调制与眼图示意}{269}
\FigureIndexRow{图24.1}{线宽增强机制示意}{280}
\FigureIndexRow{图25.1}{工程设计工作流总图}{290}
\FigureIndexRow{图25.2}{仿真工作流示意}{298}
\FigureIndexRow{图26.1}{算例桥接图}{306}
\FigureIndexRow{图27.1}{建模排错流程图}{319}
\FigureIndexRow{图28.1}{数学对象关系图}{329}
\FigureIndexRow{图29.1}{校准-验证闭环图}{347}
\FigureIndexRow{图30.1}{建模排错决策树}{361}
\FigureIndexRow{图31.1}{符号与术语分类索引图}{367}
\bottomrule
\end{longtable}
\endgroup


